% =====================================================================
%  Chapter 11: Computational Thinking and Reproducible Research
%  Source: CompMethods.tex (Overleaf) + P4 transcription (KDE) + new
% =====================================================================

\chapter{Computational Thinking and Reproducible Research}
\label{chap_Computational}

\section{The \textit{Caveat Computor} Imperative Revisited}
\label{sec_CavetComputorRevisit}

The Preface introduced the \emph{caveat computor} principle: computational
abundance makes conceptual rigour more rather than less important. This
chapter develops that principle into a practical framework for computational
regional analysis.

\noindent\textbf{What agentic coding assistants can and cannot do.}
AI coding assistants (GitHub Copilot, Claude, ChatGPT) can generate correct,
runnable code for standard analyses from plain-English descriptions. They can
suggest data sources, identify relevant packages, and debug syntax errors.
What they cannot do is determine whether the weight matrix choice is
defensible, whether the identification strategy is credible, or whether
the coefficient has a valid causal interpretation. These judgements require
the theoretical knowledge developed in Parts~I and~II.

The most dangerous analyst is not the one who cannot code---it is the one
who produces polished output without understanding its assumptions. The
obligation of this book is to ensure that the ability to generate spatial
regression code is matched by the ability to evaluate it.


\section{Reproducible Research in Practice}
\label{sec_Reproducible_ch11}

\subsection{Minimal Working Examples}

A minimal working example\index{minimal working example} (MWE) or
short, self-contained, correct example (SSCCE) is the gold standard for
communicating about code. When debugging or seeking help, an MWE:

\begin{itemize}
  \item Is \textbf{minimal}: the problem is isolated; nothing distracts from
    the error. Remove everything that is not essential to reproducing the issue.
  \item Is \textbf{working}: the code can be copied, pasted directly into a
    file, and executed without modification.
  \item Is \textbf{correct} (as far as possible): it correctly represents
    the intent of the analysis.
\end{itemize}

\noindent\textbf{Resources:}
\begin{itemize}
  \item \href{http://www.sscce.org}{sscce.org}: the authoritative reference
  \item \href{https://stackoverflow.com/help/minimal-reproducible-example}{%
    Stack Overflow guidelines}
\end{itemize}

\subsection{Literate Programming: Quarto}

Quarto (\texttt{quarto.org}) is a literate programming platform that
combines prose, code (R, Python, Julia), and mathematical notation in a
single source file, rendering to PDF, HTML, Word, or presentations. The
key workflow:

\begin{enumerate}
  \item Write code and analysis in a \texttt{.qmd} file
  \item \texttt{quarto render report.qmd} produces the output
  \item Commit \texttt{.qmd} + \texttt{renv.lock} + \texttt{requirements.txt}
    to GitHub
  \item Anyone can reproduce the analysis: \texttt{git clone} then
    \texttt{quarto render}
\end{enumerate}

\subsection{Version Control and Open Publication}

The \emph{American Economic Review} requires submission of all data and code
sufficient for replication before publication \citep{AER2024}. Urban
planning and geography journals have been slower to adopt such standards.
This book takes the position that reproducibility is a scientific norm, not
a courtesy: an analysis that cannot be replicated independently is not
science.


\section{Nonparametric Density Estimation}
\label{sec_KDE}

This section provides a technical introduction to kernel density estimation
(KDE), which arises naturally in spatial analysis (e.g.\ estimating the
density surface of events, such as crime or business locations) and in
distributional analysis (e.g.\ visualising the distribution of house prices
across neighbourhoods).

\subsection{The Histogram as a Rectangular KDE}

The histogram is the simplest nonparametric approach to density estimation.
It partitions the support of $x$ into non-overlapping bins of width $2h$
and counts observations in each bin. This is equivalent to a KDE with a
rectangular (uniform) kernel:
\begin{equation}
  K(x) = \begin{cases}
    \tfrac{1}{2} & \text{if } |x| < 1 \\
    0            & \text{otherwise}
  \end{cases}
\end{equation}
and bandwidth $h = \text{bin width}/2$.

\subsection{The General Kernel Density Estimator}

The kernel density estimator has two key ingredients:
\begin{enumerate}[label=(\roman*)]
  \item \textbf{Bandwidth} $h$: width of bin divided by 2; determines the
    size of the neighbourhood around evaluation point $x_0$.
  \item \textbf{Kernel} $K$: a weight-assigning or smoothing function
    that satisfies $\int K(x)\,dx = 1$.
\end{enumerate}

\begin{tcolorbox}[keyresult]
\begin{equation}
  \widehat{\text{pdf}}(x_0)
  = \frac{1}{Nh}\sum_{i=1}^{N} K\!\left(\frac{x_i - x_0}{h}\right)
  \label{eq_KDE}
\end{equation}
\end{tcolorbox}

\noindent Unlike the histogram, bins can \textbf{overlap}: most kernels
assign higher weight to observations near the center $x_0$ and let weights
decline with distance.

\subsection{Alternative Kernels}

\begin{figure}[htp]
\centering
\begin{tikzpicture}[>=Stealth, font=\small]
  \def\W{1.8}
  \def\H{1.0}
  \def\sep{3.8}

  \begin{scope}[xshift=0cm]
    \draw[->] (-\W-0.1,0) -- (\W+0.2,0);
    \draw[->] (0,-0.1) -- (0,\H+0.25);
    \draw[very thick] (-\W,0)--(-1,0)--(-1,\H)--(1,\H)--(1,0)--(\W,0);
    \node[below] at (0,-0.2) {Rectangular};
  \end{scope}

  \begin{scope}[xshift=\sep cm]
    \draw[->] (-\W-0.1,0) -- (\W+0.2,0);
    \draw[->] (0,-0.1) -- (0,\H+0.25);
    \draw[very thick, domain=-1:1, samples=40]
      plot (\x, {\H*(1-\x*\x)});
    \draw[very thick] (-\W,0)--(-1,0) (1,0)--(\W,0);
    \node[below] at (0,-0.2) {Epanechnikov};
  \end{scope}

  \begin{scope}[xshift=2*\sep cm]
    \draw[->] (-\W-0.1,0) -- (\W+0.2,0);
    \draw[->] (0,-0.1) -- (0,\H+0.25);
    \draw[very thick, domain=-1:1, samples=50]
      plot (\x, {\H*(1-\x*\x)^2});
    \draw[very thick] (-\W,0)--(-1,0) (1,0)--(\W,0);
    \node[below] at (0,-0.2) {Biweight};
  \end{scope}

  \begin{scope}[xshift=3*\sep cm]
    \draw[->] (-\W-0.1,0) -- (\W+0.2,0);
    \draw[->] (0,-0.1) -- (0,\H+0.25);
    \draw[very thick] (-1,0)--(0,\H)--(1,0);
    \draw[very thick] (-\W,0)--(-1,0) (1,0)--(\W,0);
    \node[below] at (0,-0.2) {Triangular};
  \end{scope}

  \begin{scope}[xshift=4*\sep cm]
    \draw[->] (-\W-0.1,0) -- (\W+0.2,0);
    \draw[->] (0,-0.1) -- (0,\H+0.25);
    \draw[very thick, domain=-1.8:1.8, samples=60]
      plot (\x, {\H*exp(-2.5*\x*\x)});
    \node[below] at (0,-0.2) {Gaussian};
  \end{scope}
\end{tikzpicture}
\caption{Five common kernel functions. All are centred at $x_0 = 0$ and
assign higher weight to nearby observations.}
\label{fig_Kernels}
\end{figure}

\subsection{The Gaussian KDE}

\begin{equation}
  \hat{f}(x) = \frac{1}{Nh}\sum_{i=1}^{N}\frac{1}{\sqrt{2\pi}}
               \exp\!\left(-0.5\left[\frac{x_i - x_0}{h}\right]^2\right)
  \label{eq_GaussianKDE}
\end{equation}

\subsection{Bandwidth Selection}

Choice of bandwidth is more important in practice than choice of kernel;
it involves a fundamental trade-off:
\begin{itemize}
  \item \textbf{Small $h$} (narrow bandwidth): low bias (correct shape)
    but high variance (jagged, imprecise estimate)
  \item \textbf{Large $h$} (wide bandwidth): low variance (smooth)
    but high bias (oversmooths the distribution)
\end{itemize}

The optimal bandwidth minimises asymptotic mean integrated squared error
(MISE). Silverman's rule-of-thumb provides a data-driven starting point:
$h = 1.06\hat\sigma n^{-1/5}$.


\section{Simulation and Resampling}
\label{sec_Simulation}

\subsection{The Bootstrap}

The bootstrap \citep{Efron1979} estimates the sampling distribution of a
statistic by resampling from the observed data with replacement. It provides
standard errors and confidence intervals without distributional assumptions,
making it particularly valuable when asymptotic approximations are unreliable.

For spatial data, the standard bootstrap must be modified to account for
spatial dependence---naive resampling destroys the spatial correlation
structure. The \textbf{spatial block bootstrap} resamples spatially
contiguous blocks of observations.

\subsection{Monte Carlo Methods for Spatial Inference}

Moran's $I$ significance is typically assessed by comparing the observed
value to a reference distribution generated by random permutation of the
data across spatial units:
\begin{enumerate}
  \item Compute observed $I^*$
  \item Randomly permute values across locations; compute $I$ for each
    permutation (e.g.\ 999 times)
  \item $p$-value = proportion of permuted $I$ values exceeding $I^*$
\end{enumerate}

This permutation approach makes no distributional assumptions and correctly
accounts for the specific topology of the weight matrix.


\section{Regional Agent-Based Models}
\label{sec_ABMs}

\subsection{What ABMs Add}

Agent-based models\index{agent-based models} (ABMs) simulate the behaviour
of individual agents (households, firms, municipalities) and allow macro-level
patterns to emerge from micro-level rules. Three phenomena are particularly
suited to ABM:

\begin{itemize}
  \item \textbf{Emergence:} Spatial patterns (segregation, clustering)
    arising from simple individual rules
  \item \textbf{Heterogeneity:} Distributional effects that aggregate models
    miss by construction
  \item \textbf{Path dependence:} Outcomes that depend on the sequence of
    historical events, not just current conditions
\end{itemize}

\subsection{The Schelling Segregation Model as a Spatial Equilibrium Problem}

\citet{Schelling1971} showed that mild individual preferences for neighbours
of the same type generate extreme aggregate segregation---a result that
anticipates the spatial equilibrium framework of Chapter~\ref{chap_SpatialEquil}
by showing how individual optimising behaviour produces spatial patterns that
no individual intended.

% TODO: Add NetLogo/nlrx worked example with calibration to ACS data


\section{Machine Learning and Causal Inference}
\label{sec_ML}

\subsection{What ML Offers: Prediction, Not Identification}

Machine learning methods (random forests, gradient boosting, LASSO, neural
networks) optimise predictive accuracy---they minimise out-of-sample
prediction error. This is valuable for several tasks in regional analysis:
variable selection, nonparametric approximation of complex response surfaces,
and out-of-sample forecasting.

What ML methods cannot do, without additional structure, is identify causal
effects. A predictive model that achieves high out-of-sample $R^2$ may do so
by exploiting spurious correlations in the training data. The fundamental
identification challenges of Chapter~\ref{chap_Causal}---selection bias,
geographic confounding, omitted variables---are not resolved by adding
flexibility to the model.

\subsection{Regularisation: LASSO and Ridge}

For high-dimensional settings with many candidate regressors, penalised
regression methods impose shrinkage on coefficients:
\begin{equation}
  \hat\beta_\text{LASSO} = \arg\min_\beta \Bigl[
    \sum(y_i - X_i'\beta)^2 + \lambda\|\beta\|_1
  \Bigr]
\end{equation}
\begin{equation}
  \hat\beta_\text{Ridge} = \arg\min_\beta \Bigl[
    \sum(y_i - X_i'\beta)^2 + \lambda\|\beta\|_2^2
  \Bigr]
\end{equation}

LASSO ($L_1$ penalty) produces sparse solutions (many coefficients set
exactly to zero) and serves as a variable selection tool. Ridge ($L_2$
penalty) shrinks all coefficients continuously toward zero without setting
any to exactly zero.

\subsection{Causal Forests and Heterogeneous Treatment Effects}

\citet{Wager2018} combine the random forest algorithm with the
potential-outcomes framework to estimate
\textbf{heterogeneous treatment effects}\index{heterogeneous treatment
effects}---the individual-level causal effect $\tau_i = Y_{1i} - Y_{0i}$---
as a function of covariates. The causal forest modifies the splitting rule of
a random forest to maximise treatment effect heterogeneity rather than
predictive accuracy, and uses orthogonalisation (the Robinson transformation)
to partial out confounders. This connects directly to the GWR framework:
where GWR asks whether the \emph{association} between $X$ and $Y$ varies
across space, a spatial causal forest asks whether the \emph{causal effect}
of a treatment varies across space.

\renewcommand\bibname{Further reading}
\begin{subappendices}
\bibliographystyle{econometrica}
\bibliography{Literature/OverLord}
\IfFileExists{Literature/CompMethods.bbl}{\chapter{Computational Methods}\label{chap_CompMethods}

\section{Principles of coding}
\subsection{Debugging code}
Most questions about code are made more answerable by the addition of a minimal working example (MWE)\index{MWE|see{minimal working example}}\index{minimal working example} or a short, self-contained, correct example (SSCCE).

Your debugging question may have an obvious answer, or it may not. When enlisting the help of others in debugging\index{code!debugging} your code, think about it from the point of view of the person whose help you're seeking. To ensure that the answer she provides is sound, she will probably want to create a sample document, make the changes she recommends, and check that it works. If she thinks she knows the answer but doesn't have the time to create a sample document, she may in the interest of quality control refrain from answering at this time. Thus you should always provide a block of code that can be copied, pasted directly into a file, compiled, tweaked, corrected, and pasted back.

In this context, ``minimal'' means that the problem is isolated and there is nothing in the sample that distracts from the error you have or the effect you desire. ``Working'' doesn't mean the problem is solved, it means the code sample can be copied-and-pasted and directly compiled. The site \href{http://www.sscce.org}{sscce.org} is dedicated to the topic. There are also several articles on the topic as it applies specifically to LaTeX:

\begin{itemize}
\item \href{http://theoval.cmp.uea.ac.uk/~nlct/latex/minexample/minexample.pdf}{``Creating a LaTeX Minimal Example''} (PDF) from Dr. Nicola Talbot's LaTeX resources,
\item \href{http://www.tex.ac.uk/cgi-bin/texfaq2html?label=minxampl}{``How to make a `minimum example'"} from the UK \TeX{} FAQ,
\item \href{http://www.minimalbeispiel.de/mini-en.html}{``What is a minimal working example?"} from the de.comp.text.tex newsgroup.
\end{itemize}


\section{High-performance computing}
This section introduces the concept of high-performance computing\index{high-performance computing} and illustrates an application using the University of Michigan's high-performance Linux cluster Flux\index{Flux|see{software}}\index{Flux|seealso{high-performance computing}}\index{software!Flux}.


\subsection{Parallelised statistical computations}
Many modern statistical techniques rely on computations that are parallel. Resampling methods (such as cross-validation and the bootstrap) and many simulations can be separated into independent sets of computations. 

As an example, we use simulation to produce a power curve for an independent samples $t$-test. After demonstrating the simulation using the \texttt{lapply} function in \textsf{R}, we parallelise the simulation using the \texttt{mcapply} function from the \texttt{multicore} package.



\subsection{University of Michigan high-performance cluster}
\subsubsection{Logging into the cluster and copying files}
\subsubsection{Software on Flux}
\subsubsection{Description of key commands}
\subsubsection{Running a batch job}

\section{Nonparametric econometrics}
This section with provide a basic introduction to nonparametric techniques in econometrics.

\linenumbers
\footnotesize
\begin{verbatim}
R> y <- plot(density(rnorm))
R> lines(x, add=TRUE)
\end{verbatim}
\normalsize
\nolinenumbers
}{\typeout{Need bibtex on CompMethods}}
\end{subappendices}
